\documentclass[../../main.tex]{subfiles}
\graphicspath{{../../image/}{../../figure/}}
\begin{document}

\section{波扇容许性准则}\label{section:ch09-波扇容许性准则}

关于容许性的各个方面,对一般弱解的讨论已见于\cref{chapter:ch04-cauchy问题},对单个激波的讨论已见于\cref{chapter:ch08-容许激波}。因而我们已经遇到了一系列容许性准则,并看到它们紧密相关但并不完全等价。正如我们后文将看到的,这些准则中最具判别力者,即粘性激波剖面与 Liu $E$ 条件,只要我们局限于严格双曲方程组与中等强度激波,就足以剔除一切赝解。然而一旦转到非严格双曲方程组和/或带强激波的解,局面就变得模糊不清。容许性问题至今仍未完全解决。

任何理性的新容许准则都应恪守某些基本原则,这些是积累经验的结晶。它们包括:

\subsection{局部化}\label{subsection:ch09-局部化}
解的容许性检验应单独施加于定义域中每一点 $(\bar{x},\bar{t})$,因而只涉及解在 $(\bar{x},\bar{t})$ 的任意小邻域上的限制,例如圆 $\{(x,t):|x-\bar{x}|^{2}+|t-\bar{t}|^{2}<r^{2}\}$,其中 $r$ 固定但可任意小。这与双曲方程组解应具有局部依赖性质这一一般原理相容。

\subsection{演化性}\label{subsection:ch09-演化性}
容许性检验应面向未来,不顾及过去。因而解在点 $(\bar{x},\bar{t})$ 处的容许性应仅取决于它在半圆盘 $\{(x,t):|x-\bar{x}|^{2}+|t-\bar{t}|^{2}<r^{2},\ t\geqslant\bar{t}\}$ 上的限制。这与时间不可逆性原理一致,该原理贯穿了我们迄今遇到的所有容许性准则,如熵、粘性等。

\subsection{平移不变性}\label{subsection:ch09-平移不变性}
解 $U$ 在 $(\bar{x},\bar{t})$ 处容许,当且仅当平移后的解 $\bar{U}$,即 $\bar{U}(x,t)=U(x+\bar{x},t+\bar{t})$,在原点 $(0,0)$ 处容许。

\subsection{伸缩不变性}\label{subsection:ch09-伸缩不变性}
解 $U$ 在 $(0,0)$ 处容许,当且仅当对每个 $\alpha>0$,伸缩后的解 $\bar{U}_{\alpha}$,即 $\bar{U}_{\alpha}(x,t)=U(\alpha x,\alpha t)$,在 $(0,0)$ 处容许。

我们把注意力集中于这样一类弱解 $U$:对定义域中每一固定点 $(\bar{x},\bar{t})$,极限
\begin{align}
\bar{U}(x,t)=\lim_{\alpha\downarrow0}U(\bar{x}+\alpha x,\bar{t}+\alpha t)
\label{align:ch09-9-2-1}
\end{align}
在 $L^{1}_{\mathrm{loc}}((-\infty,\infty)\times(0,\infty))$ 中存在。注意此时 $\bar{U}$ 必为\eqref{align:ch09-9-1-1}的一个自相似解。本着上述原则的精神,可用 $\bar{U}$ 在原点的容许性作为 $U$ 在点 $(\bar{x},\bar{t})$ 处容许性的检验。由于 $\bar{U}$ 描绘了一簇从原点辐射的波,这类检验就构成\emph{波扇容许性准则}。

在\eqref{align:ch09-9-2-1}中取极限,无异于用一架聚焦于点 $(\bar{x},\bar{t})$ 的显微镜观察解 $U$。当 $U\in BV_{\mathrm{loc}}$ 时,由\cref{theorem:ch01-BV链式法则},该极限存在,可能的例外是在一维 Hausdorff 测度为零的不规则点集 $(\bar{x},\bar{t})$ 上。特别地,若 $(\bar{x},\bar{t})$ 是 $U$ 的近似连续点,则 $\bar{U}(x,t)$ 为常值状态 $U_{0}$;若 $(\bar{x},\bar{t})$ 是 $U$ 的近似跳跃间断点,则当 $x<st$ 时 $\bar{U}(x,t)=U_{-}$,当 $x>st$ 时 $\bar{U}(x,t)=U_{+}$,其中 $U_{\pm}$ 是 $U$ 的近似单侧极限,而 $s$ 是 $(\bar{x},\bar{t})$ 处跳跃间断的斜率。至于在 $U$ 的不规则点 $(\bar{x},\bar{t})$ 处该极限是否也存在、而所得波扇 $\bar{U}$ 应更复杂,将在后文研究:对真正非线性的标量守恒律在\cref{chapter:ch11-一维标量守恒律},对真正非线性的两个守恒律方程组在\cref{chapter:ch12-两方程真非线性组},对一般真正非线性守恒律方程组在\cref{chapter:ch14-前沿追踪与标准riemann半群}。

正如\cref{section:ch09-自相似解与riemann问题}中所见,波扇 $\bar{U}$ 一般是常值状态、激波与中心稀疏波的复合体。最简单的波扇容许准则假定:波扇容许,当且仅当其每个激波单独满足\cref{chapter:ch08-容许激波}中讨论的激波容许条件。正如我们在下一节将看到的,这在许多情形下是充分的。另一些把波扇视为一个整体而非单个波之集合的波扇容许准则,包括熵率条件与粘性波扇剖面检验,将在后文讨论。

\end{document}
